\section{Discussão}
\label{sec:discussao}

% Na discussão, analise os resultados obtidos, comparando-os com estudos
% anteriores, destacando a contribuição do trabalho e suas implicações. Também
% pode abordar limitações e sugestões para trabalhos futuros.

\subsection{Limitações}
\label{subsec:implementation_limitations}

A principal limitação observada está relacionada ao custo computacional das
interações entre partículas, que aumenta conforme a quantidade de partículas é
ampliada. Esse comportamento limita a escalabilidade da implementação e reduz a
taxa de quadros para configurações com maior número de partículas.

Outra limitação está relacionada ao custo adicional da renderização. Embora a
versão \textit{headless} apresente desempenho superior, a execução completa
precisa realizar tanto os cálculos da simulação quanto a representação gráfica
das partículas.

O programa também utiliza uma implementação falha de multi-thread. O que
explica os picos negativos em fps e tempo de passo, parte da simulação tem
problemas tentando acessar dados em outras threads resultando em perda abruta
de dados de algumas partículas.

Além disso, não foram incorporadas técnicas de otimização presentes em
simuladores mais avançados, como aceleração em GPU, estruturas espaciais
especializadas para busca de vizinhos e refinamento adaptativo. A utilização
dessas técnicas poderia reduzir o custo das interações e permitir a execução de
simulações com maior quantidade de partículas mantendo uma taxa de quadros mais
elevada.
